\section{Introduction}\label{sec:intro}

Kinetic reduced magnetohydrodynamics (KRMHD), the limit $k_\perp\rho_i\to0$ of gyrokinetics derived by \citet{Schekochihin2009}, hereafter S09,\footnote{Appendix~\ref{app:notation} compares our notation with that of S09 and of \citet{Chandran2026}.} describes anisotropic fluctuations at scales much larger than the ion gyroradius $\rho_i$ in a strongly magnetised, weakly collisional plasma. Its Alfv\'enic fluctuations obey reduced magnetohydrodynamics (RMHD). Its compressive fluctuations, the slow waves and the entropy mode, obey a kinetic equation for the ion distribution in the parallel velocity $v_\parallel$; they are carried by the Alfv\'enic fluctuations and do not act back on them. Without collisions they are damped by phase mixing \citetext{\citealp{Barnes1966}; S09, \S6.2}, which moves their free energy to ever finer structure in $v_\parallel$, that is, to large Hermite order $m$ in a Hermite expansion in $v_\parallel$ \citep{ZoccoSchekochihin2011,Kanekar2015,Schekochihin2016}. In turbulence the same free energy is also carried to small perpendicular scales. Which of the two routes it takes decides whether the compressive cascade heats ions through Landau damping or reaches the ion gyroscale \citep[S09;][]{Schekochihin2016,Meyrand2019,Kawazura2019,Kawazura2020}, and the Hermite spectrum $W(m)$ is the most direct record of that competition.

Two pictures of $W(m)$ in turbulence have been proposed. In the first, nonlinear advection couples fluctuations with parallel wavenumbers of opposite sign and so returns free energy from high to low $m$, as a stochastic plasma echo \citep{KanekarThesis2015,Schekochihin2016,Parker2016,AdkinsSchekochihin2018}. \citet{Meyrand2019} found in KRMHD turbulence at $\beta_i=1$ a spectrum somewhat steeper than $m^{-1}$, against $m^{-1/2}$ for linear phase mixing, and a net flux to high $m$ much smaller than the injection rate away from the forcing scales, and attributed both to such echoes. In gyrokinetic simulations \citet{Kawazura2019,Kawazura2020} found $W(m)\propto m^{-1/2}$ at high $\beta_i$ and close to $m^{-1}$ at $\beta_i=0.1$, and read the steeper spectrum as phase mixing suppressed by echoes. In the second picture (S09, \S6.3), the compressive fluctuations have no parallel cascade: the field-aligned wavenumber at which they were injected stays fixed while the Alfv\'en waves carry them to smaller perpendicular scales, so at small enough scales the perpendicular cascade outruns the damping. The two pictures differ on whether the compressive fluctuations acquire structure along the moving field lines. \citet{Meyrand2019} argued that they do, because resistive reconnection relabels the field lines within about one eddy-turnover time and the compressive fluctuations then inherit the parallel cascade of the Alfv\'enic ones.

We make the second picture quantitative and test it. Section~\ref{sec:slices} shows that the compressive sector of KRMHD is exactly a continuum of passive scalars, or slices, labelled by the parallel phase velocity $u$, each carried by the $\boldsymbol E\times\boldsymbol B$ flow and streaming along the perturbed field at its own speed. The nonlinearity therefore moves no free energy between phase velocities, and free energy moves between Hermite moments only as the correlations between slices change (section~\ref{sec:hermite}). When the induction equation is ideal, the field lines are frozen into the flow that carries all the slices, and every volume-integrated single-point statistic, including $W(m)$, evolves as in linear phase mixing along the initial field lines (section~\ref{sec:frozen}). In turbulence the slices are also cascaded to small perpendicular scales and dissipated there. Section~\ref{sec:twochannel} proposes that this is the only effect of the nonlinearity on $W(m)$: every Hermite moment loses free energy at the same rate $1/T$, the rate at which the turbulence carries a passive scalar to the dissipation scale, while phase mixing proceeds at the field-aligned wavenumber $\kappa$ of the injected free energy. The steady state is
\begin{equation}\label{eq:intro-W}
  W(m)=\frac{\eps}{s\kappa\sqrt{2m}}\exp\!\Big(-\frac{\sqrt{2m}}{s\kappa T}\Big),\qquad m_*=\frac{(s\kappa T)^2}{2}\sim\frac{\beta_i}{\chi^2},
\end{equation}
for injection rate $\eps$, $s=\vth/v_A=\beta_i^{1/2}$ and nonlinearity parameter $\chi$, summed over the distribution of $\kappa$. Two times set the spectrum: free energy injected at $\kappa$ phase mixes up to moment $m$ in the time $\sqrt{2m}/(s\kappa)$, which is shorter at higher $\beta_i$, and the perpendicular cascade removes it in the time $T$, which is set by the Alfv\'enic turbulence. The exponential in (\ref{eq:intro-W}) is the fraction of the free energy that survives the sink while it phase mixes up to $m$, and $m_*$ is the moment it reaches in the time $T$. Every quantity in (\ref{eq:intro-W}) can be measured independently of $W(m)$, so the model has no free parameter. It predicts that the perpendicular dissipation of each Hermite moment is proportional to $W(m)$, that $W(m)$ and the Hermite flux computed along the field follow the model with no deficit at large $m$, that at fixed $\chi$ the break $m_*$ grows with $\beta_i$, in proportion to it if $T$ does not depend on $\beta_i$, and that the field-aligned spectrum of the free energy stays that of the injection as the resolution increases.

Section~\ref{sec:simulations} tests these predictions in forced KRMHD turbulence with the GANDALF code \citep{Kanekar2025}. At $\beta_i=1$ and $\chi\approx1$ the perpendicular dissipation per Hermite moment is proportional to $W(m)$ to within 20\%, and the model reproduces $W(m)$ to within 8\% for $m\le30$ and the field-aligned flux to within 7\% for $m\le24$. At $\beta_i=1$ the break $m_*$ lies beyond the $M=32$ moments we resolve. At $\beta_i=0.3$ and $0.1$, with the same Alfv\'enic turbulence, it moves into the resolved range as predicted, and $W(m)$ steepens accordingly. There the measured spectrum lies above the model at every $m$, whereas a return flux would put it below the model at large $m$. The excess comes from the assumption of a single time $T$. The sink switches on over about five outer eddy-turnover times, as the cascade carries fresh free energy to the dissipation scale, so it removes less of the free energy at small $m$, which phase mixing brought there more recently. At lower $\beta_i$ phase mixing is slower and the free energy at a given $m$ is older, so the sink differs between the runs at a given $m$ but is the same function of the time since injection in all of them. Doubling the number of Hermite moments leaves $W(m)$ unchanged where both resolutions resolve it. Raising the perpendicular resolution does not: up to $256^3$ grid points $T$ grows by about the same amount with each doubling, so the sink is not independent of the dissipation at the resolutions we reach. The other approximation of the model, a frozen field-aligned wavenumber, holds at every resolution: at $\beta_i=0.3$ the field-aligned spectrum of the free energy does not change from $64^3$ to $256^3$, apart from a tail holding less than 0.5\% of it. With the $T$ measured at each resolution the model still describes $W(m)$, and there is still no return flux. With the parallel flow forced in place of the density, as by \citet{Meyrand2019}, the spectrum steepens at small $m$, by the amount the model gives with the injected distribution of $\kappa$. With the Hermite resolution raised until the break is resolved, $m_*$ grows as $\beta_i^{1.09\pm0.04}$ for $0.05\le\beta_i\le0.3$ at $\chi\approx1$. At $\beta_i=0.05$ it grows somewhat faster than $\chi^{-2}$ down to $\chi\approx0.56$ and about as $\chi^{-4}$ between $\chi=0.56$ and 0.40. In units of the outer eddy time the sink weakens at the smallest $\chi$, so the estimate $m_*\sim\beta_i/\chi^2$, which takes $T$ proportional to the eddy time, holds only near $\chi\approx1$.

The ingredients of (\ref{eq:intro-W}) are not new, and we credit them where they are used. The rate comparison between Landau damping at a frozen parallel wavenumber and the perpendicular cascade is S09's (\S6.3, eqs~217--220). The functional form is that of \citet[eq.~112]{ZoccoSchekochihin2011} for growing modes, with their growth rate replaced by the cascade rate, and \citet[eqs~2.58--2.59]{Schekochihin2016} solved the same continuity equation in $m$ with a collisional sink. The break scale is the threshold $k_\parallel\vth\sim\sqrt m\,k_\perp u_\perp$ of \citet[eq.~4.5]{Schekochihin2016}, and \citet{Kawazura2019} stated $m\sim\beta_i$ for critically balanced turbulence. A Hermite spectrum decaying exponentially in $\sqrt m$, at a rate proportional to the ratio of the cascade rate to $k_\parallel\vth$, was seen in simulations of a kinetic passive scalar advected by a two-dimensional flow \citep[chapter~4]{KanekarThesis2015}. What is new here is the exact statement that in three-dimensional KRMHD with ideal induction the nonlinearity moves no free energy in $m$, the model with independently measured inputs that follows from it, and its test in turbulence with diagnostics computed along the field. The slice decomposition itself has precedents in the G-transform of \citet{MorrisonPfirsch1992} and \citet{Morrison2000}, its application to drift kinetics by \citet{Heninger2018}, and the per-velocity invariants of \citet{PlunkParker2014} and \citet{Plunk2015}.

Section~\ref{sec:model} sets out the compressive sector of KRMHD, sections~\ref{sec:slices} to~\ref{sec:frozen} give the exact results, section~\ref{sec:twochannel} the model and its predictions, section~\ref{sec:simulations} the simulations, and section~\ref{sec:discussion} discusses what the results do and do not establish.
